\subsection{P014 --- CO2- and PM2.5-Focused Optimal Ventilation Strategy Based on Predictive Control}
\label{paper:P014}
\textbf{Authors:} Young Jae Choi, Eun Ji Choi, Jae Yoon Byun, Hyeun Jun Moon, Min Ki Sung, Jin Woo Moon\par
\textbf{Major Category:} Building Environment and IEQ\par
\textbf{Secondary Category:} HVAC and Building Services, AI and Machine Learning for Buildings\par
\textbf{Keywords:} CO2, PM2.5, Ventilation, Predictive control, VAV, Indoor air quality, Co-simulation\par
\paragraph{主な主張}
occupancy-based system-level DCVと, DNNによるCO2/PM2.5 one-step predictionを用いたzone-level controlを組み合わせたVAV ventilation strategyを提案した. EnergyPlus-CONTAM-Python co-simulationで, CO2は100\%, PM2.5は97.33\%の時間で上限値以下を維持した. \cite{Choi:2025CO2PM25Ventilation}
\paragraph{新規性}
CO2とPM2.5の予測値と上限値のerror rateからcontrol priorityを動的に切替え, 複雑なoptimization functionへの依存を抑えたmulti-pollutant ventilation controlを構成した点に特徴がある.
\paragraph{位置付け}
CO2とPM2.5を同時に扱うIAQ controlが主要目的であるためBuilding Environment and IEQを主分類とし, ventilation/HVACとpredictionを副分類とする.
